\chapter{Proper Purity and the Decomposition Theorem}
\label{chap:decomposition}

The arithmetic estimate is now complete:
\cref{thm:deligne-laumon-direct-image} proves that \(Rf_!\) preserves upper
weights.  We may therefore turn the formal perverse machinery into purity,
intersection-cohomology statements, and decomposition.  This order is the
logical heart of the notes.

\section{The direct-image theorem in perverse form}

\begin{corollary}[Perverse upper direct image]
\label{cor:perverse-upper-direct-image}
Let \(f_0:X_0\to Y_0\) be separated of finite type.  If
\(\cP_0\in\Perv(X_0)\) has weights at most \(w\), then
\[
  \pH^a(Rf_{0!}\cP_0)
  \quad\text{has weights at most }w+a
\tag{16.1}\label{eq:perverse-rfshriek-upper}
\]
for every \(a\).
\end{corollary}

\begin{proof}
The derived statement of
\cref{thm:deligne-laumon-direct-image} says
\(Rf_{0!}\cP_0\in D_{\leq w}(Y_0)\).  Apply the perverse weight criterion
\cref{thm:perverse-weight-criterion}.
\end{proof}

\begin{corollary}[Dual lower direct image]
\label{cor:perverse-lower-direct-image}
If \(\cK_0\in D_{\geq w}(X_0)\), then
\[
  Rf_{0*}\cK_0\in D_{\geq w}(Y_0).
\tag{16.2}\label{eq:rfstar-lower}
\]
Equivalently, if a perverse sheaf \(\cP_0\) has weights at least \(w\), then
\[
  \pH^a(Rf_{0*}\cP_0)
  \quad\text{has weights at least }w+a.
\tag{16.3}\label{eq:perverse-rfstar-lower}
\]
\end{corollary}

\begin{proof}
By definition \(\D_{X_0}\cK_0\in D_{\leq-w}\).  Apply the upper theorem and
the Verdier-duality isomorphism
\[
  \D_{Y_0}Rf_{0*}\cK_0
  \simeq Rf_{0!}\D_{X_0}\cK_0.
\]
The right side has weights at most \(-w\); dualizing is exactly
\eqref{eq:rfstar-lower}.  Formula
\eqref{eq:perverse-rfstar-lower} follows from
\cref{thm:perverse-weight-criterion}.
\end{proof}

\begin{theorem}[Proper direct-image purity]
\label{thm:proper-direct-image-purity}
Let \(f_0:X_0\to Y_0\) be proper and let \(\cK_0\) be pure of weight \(w\).
Then \(Rf_{0*}\cK_0\) is pure of weight \(w\).  In perverse degrees,
\[
  \pH^a(Rf_{0*}\cK_0)
  \quad\text{is pure of weight }w+a.
\tag{16.4}\label{eq:proper-perverse-purity}
\]
\end{theorem}

\begin{proof}
Properness gives \(Rf_{0!}=Rf_{0*}\).  The upper direct-image theorem gives
weights at most \(w\); \cref{cor:perverse-lower-direct-image} gives weights at
least \(w\).  Their intersection is purity.  Apply
\eqref{eq:pH-pure-shift} for the last assertion.
\end{proof}

\begin{perversebox}
Upper weights come from the Fourier/curve argument.  Lower weights come from
Verdier duality.  Properness is the point at which the two functors become the
same and the two inequalities meet.  Merely knowing the upper theorem never
implies purity for a nonproper map.
\end{perversebox}

\section{Purity of intermediate extensions}

\begin{theorem}[Middle extension preserves purity]
\label{thm:jmiddle-purity}
Let \(j_0:U_0\hookrightarrow X_0\) be an open immersion and let
\(\cP_0\in\Perv(U_0)\) be pure of weight \(w\).  Then
\[
  j_{0!*}\cP_0
  \quad\text{is pure of weight }w.
\tag{16.5}\label{eq:jmiddle-pure}
\]
\end{theorem}

\begin{proof}
The upper direct-image theorem applied to \(j_0\) shows that
\(j_{0!}\cP_0\) has weights at most \(w\).  Hence
\({}^pj_{0!}\cP_0=\pH^0(j_{0!}\cP_0)\) has weights at most \(w\).
Duality and the same theorem show that
\({}^pj_{0*}\cP_0\) has weights at least \(w\).

By definition \(j_{0!*}\cP_0\) is a quotient of
\({}^pj_{0!}\cP_0\) and a subobject of
\({}^pj_{0*}\cP_0\).  Upper bounds pass to quotients and lower bounds pass to
subobjects.  Thus it has both bounds and is pure.
\end{proof}

\begin{corollary}[Purity of intersection complexes]
\label{cor:IC-purity}
Let \(X_0\) be irreducible, let \(j_0:U_0\hookrightarrow X_0\) be a smooth
dense open of dimension \(d\), and let \(\cL_0\) be lisse and punctually pure
of weight \(u\).  Then
\[
  \IC_{X_0}(\cL_0)=j_{0!*}(\cL_0[d])
  \quad\text{is pure of weight }u+d.
\tag{16.6}\label{eq:IC-pure-weight}
\]
\end{corollary}

\begin{proof}
Use \cref{prop:lisse-perverse-weight-shift,thm:jmiddle-purity}.
\end{proof}

\section{Pure intersection cohomology}

\begin{corollary}[Proper intersection-cohomology purity]
\label{cor:proper-IH-purity}
Let \(X_0\) be proper of pure dimension \(d\), and let \(\cL_0\) be as in
\cref{cor:IC-purity}.  Then
\[
  IH^r(X,\cL)
  :=\mathbb H^r(X,\IC_X(\cL))
\]
is pure of weight
\[
  u+d+r.
\tag{16.7}\label{eq:IH-pure}
\]
\end{corollary}

\begin{proof}
Apply \cref{thm:proper-direct-image-purity} to the structural morphism
\(a_0:X_0\to\Spec\Fq\) and the pure complex
\(\IC_{X_0}(\cL_0)\) of weight \(u+d\).  On a point, perverse and ordinary
cohomology agree; \eqref{eq:point-weight-shift} gives
\eqref{eq:IH-pure}.
\end{proof}

\begin{corollary}[Smooth-proper cohomological Riemann hypothesis]
\label{cor:perverse-smooth-proper-purity}
If \(X_0/\Fq\) is smooth and proper of pure dimension \(d\), then
\[
  H^i(X,\Qbarl)\quad\text{is pure of weight }i
\tag{16.8}\label{eq:perverse-smooth-purity}
\]
for every \(i\).
\end{corollary}

\begin{proof}
Smoothness gives \(\IC_{X_0}=\Qbarl[d]\), pure of weight \(d\).  Apply
\cref{thm:proper-direct-image-purity} to \(a_0:X_0\to\Spec\Fq\):
\[
  Ra_{0*}\Qbarl[d]\quad\text{is pure of weight }d.
\]
Its degree-\(r\) cohomology is
\[
  H^r(Ra_{0*}\Qbarl[d])=H^{r+d}(X,\Qbarl)
\]
and has weight \(d+r\).  Put \(i=d+r\).
\end{proof}

\section{Geometric semisimplicity}

\begin{theorem}[BBD geometric semisimplicity]
\label{thm:pure-perverse-geometric-semisimple}
Let \(\cP_0\) be a pure perverse Weil sheaf on \(X_0\).  After base change to
\(\Fbar\) and forgetting the Frobenius structure, \(\cP\) is a finite direct
sum of simple perverse sheaves:
\[
  \cP\simeq
  \bigoplus_\lambda
  \IC_{\overline{S_\lambda}}(\cL_\lambda)^{\oplus m_\lambda}.
\tag{16.9}\label{eq:pure-perverse-semisimple}
\]
\end{theorem}

\begin{proof}
Finite length is \cref{thm:perverse-finite-length}.  Let \(\cP'\subset\cP\)
be a maximal semisimple geometric subobject and put
\(\cP''=\cP/\cP'\).  After replacing the ground field by a finite extension,
the object, subobject, quotient, and exact sequence all descend:
\[
 0\longrightarrow\cP'_0\longrightarrow\cP_0
 \longrightarrow\cP''_0\longrightarrow0.
\]
Because \(\cP_0\) is pure of weight \(w\), the quotient \(\cP''_0\) has
weights \(\leq w\), while the subobject \(\cP'_0\) has weights \(\geq w\).
The sequence defines
\[
 e_0\in\Ext^1_{X_0}(\cP''_0,\cP'_0).
\]
BBD 5.1.15(iii) says that, for a first argument of weights \(\leq w\) and a
second argument of weights \(\geq w\), every positive-degree arithmetic
extension maps to zero after geometric base change.  Hence the image of
\(e_0\) in \(\Ext^1_X(\cP'',\cP')\) vanishes and the geometric sequence
splits.  If \(\cP''\neq0\), a simple subobject of the finite-length object
\(\cP''\) would enlarge \(\cP'\), contradicting maximality.  Thus
\(\cP''=0\).  The simple-object classification
\cref{thm:simple-perverse-classification} now gives
\eqref{eq:pure-perverse-semisimple}.  This is BBD 5.3.8.
\end{proof}

\begin{warning}[Geometric is not Frobenius semisimplicity]
\label{warn:frobenius-not-semisimple}
The theorem forgets the Weil structure.  It does not say that Frobenius acts
semisimply.  On a point, take a two-dimensional representation on which
Frobenius is \(\alpha\begin{psmallmatrix}1&1\\0&1\end{psmallmatrix}\), with
\(\alpha\) pure.  All eigenvalues have the same weight, but the Frobenius
operator has a nontrivial Jordan block.
\end{warning}

\section{Splitting a pure complex}

\begin{theorem}[BBD splitting of pure complexes]
\label{thm:pure-complex-splitting}
If \(\cK_0\in\Dbc(X_0,\Qbarl)\) is pure, then after base change to \(\Fbar\)
there exists an isomorphism
\[
  \cK\simeq\bigoplus_a\pH^a(\cK)[-a].
\tag{16.10}\label{eq:pure-complex-splitting}
\]
The objects \(\pH^a(\cK)\) are canonical; the splitting generally is not.
\end{theorem}

\begin{proof}
Assume \(\cK_0\) is pure of weight \(w\), and let \(a\) be its largest
nonzero perverse degree.  Put
\[
 A_0=\pH^a(\cK_0)[-a],
 \qquad
 B_0=\ptau_{\leq a-1}\cK_0.
\]
The perverse weight criterion makes \(A_0\) pure of weight \(w\).  It makes
\(B_0\) pure of weight \(w\) as well, since its perverse cohomology objects
are precisely \(\pH^r(\cK_0)\), pure of weight \(w+r\), for \(r<a\).
The truncation triangle is
\[
 B_0\longrightarrow\cK_0\longrightarrow A_0
 \xrightarrow{\partial_0}B_0[1].
\tag{16.11}\label{eq:perverse-truncation-obstruction}
\]
Its connecting morphism is an arithmetic extension class
\[
 [\partial_0]\in\Ext^1_{X_0}(A_0,B_0).
\]
Apply BBD 5.1.15(iii) with \(A_0\) of weights \(\leq w\) and \(B_0\) of
weights \(\geq w\): the image of this class in
\(\Ext^1_X(A,B)\) after geometric base change is zero.  Therefore the
geometric truncation triangle splits.  Induction on the finite perverse
amplitude of \(\cK_0\) gives \eqref{eq:pure-complex-splitting}.  This is
BBD 5.4.5.
\end{proof}

\begin{theorem}[Decomposition theorem over a finite field]
\label{thm:decomposition-finite-field}
Let \(f_0:X_0\to Y_0\) be proper and let \(\cK_0\) be pure.  After base change
to \(\Fbar\),
\[
  Rf_*\cK\simeq
  \bigoplus_a \pH^a(Rf_*\cK)[-a],
\tag{16.12}\label{eq:decomposition-theorem}
\]
and every perverse cohomology object on the right is semisimple.
In particular, for \(X_0\) pure dimensional and \(\cL_0\) pure on a smooth
dense open, this applies to \(\cK_0=\IC_{X_0}(\cL_0)\).
\end{theorem}

\begin{proof}
Proper direct-image purity is
\cref{thm:proper-direct-image-purity}.  Apply
\cref{thm:pure-complex-splitting} to \(Rf_{0*}\cK_0\), and
\cref{thm:pure-perverse-geometric-semisimple} to each perverse cohomology
object.
\end{proof}

\begin{remark}[Relative hard Lefschetz]
For a projective morphism and an ample class \(\eta\), the stronger theorem
gives
\[
  \eta^r:\pH^{-r}(Rf_*\IC_X)
  \xrightarrow{\sim}\pH^r(Rf_*\IC_X)(r).
\]
Together with primitive decomposition it makes the splitting more structured.
It is not needed for the classical Weil conjectures and is not used anywhere
upstream.  Keeping it outside the dependency chain avoids concealing an
additional deep theorem behind the word “decomposition.”
\end{remark}

\begin{formalization}
The theorem interfaces should be layered:
\[
\begin{array}{c}
 Rf_!\text{ upper weights}\\
 \Downarrow\ \text{duality}\\
 Rf_*\text{ lower weights}\\
 \Downarrow\ \text{properness}\\
 Rf_*\text{ purity}\\
 \Downarrow\\
 \text{IC/IH purity, geometric semisimplicity, decomposition.}
\end{array}
\]
The decomposition is an existence of an isomorphism after base change, not a
chosen definitional equality and not a Frobenius-equivariant semisimplicity
record.  The smooth test must reduce
\(Ra_*\Qbarl[d]\) to \(H^{r+d}\) with weight \(d+r\).
\end{formalization}
